\section{Purpose and model}\label{sec:introduction}
The first paper's Theorem P$'$ bounds the information delivered by an adaptive
machine using few monotone passes~\cite{BrendeTranspose}. Its explicit
inequality includes the runtime $T$. The paper obtains
$r\ge\log_{k+1}d-O_M(1)$ when $\log T=O(b)$ and identifies removal of the
runtime premise as unfinished formal work. This memo supplies that step for
halting runs. It leaves the original finite-time and nonhalting-prefix
statements intact.

Bounds exponential in the number of reversals are standard. In particular,
Hernich and Schweikardt state a run-length bound for reversal- and
space-bounded machines in Lemma~2.3, attributing the deterministic argument
to earlier work~\cite{HernichSchweikardt2006}. Our proof works directly with
the first paper's deterministic machine, including simultaneous head moves,
stationary heads, and unreadable output. We make no priority claim for the
run-length principle.

The model has $q+1$ states, a read-only binary input tape, $k$ initially blank
work tapes with $g+1$ symbols, and an append-only output that the machine
cannot read. Each input or work head moves by $-1$, $0$, or $1$ per step;
all start at position zero. A pass is an interval on which each of these
$k+1$ head trajectories is monotone, with a separate direction for each
head. Empty passes are allowed. A run halting at $T$ has a transition at
every step before $T$ and none at $T$.

We use $d$ for the matrix dimension, $b$ for its entry width, and $L=d^2b$.
The canonical width is $b(d)=\max(1,\lceil\log_2d\rceil)$, as in the frozen
statement of the first paper. For one fixed machine define
\begin{align}
 S_M&=3(q+1)(g+1)^k,& C_M&=(2k+3)(S_M+1),\label{eq:machine-constants}\\
 \sigma_M&=\lceil\log_2(q+1)\rceil+k\lceil\log_2(g+1)\rceil,&
 \ell_M&=\lceil\log_2 C_M\rceil,\nonumber\\
 A_M&=2+\sigma_M+(k+4)(5+\ell_M).\label{eq:pass-constant}
\end{align}
The Lean declarations are authoritative proofs; the outlines below explain
them. This memo imports the original library by a relative path and does
not alter its frozen machine definitions.

\section{Native run length}\label{sec:run-length}

\begin{lemma}[Clean-ray intervals]\label{lem:ray}
\formal{RaySimulation.lean}{ray_simulation, ray_eventually_none, ray_interval_bound}
At the start of an interval, assign each head either the instruction to stay
at its current cell or a direction along a ray. For a moving input head,
assume the whole ray lies outside the input word. For a moving work head,
assume every cell strictly ahead on the ray is blank. During the interval,
allow each head only stationary moves or moves along its assigned ray.
If the actual run eventually halts, any such interval of executed steps has
length at most $S_M$.
\end{lemma}
\begin{outline}
The signature consists of the state, the currently scanned input symbol
(including the outside-word symbol), and the $k$ scanned work symbols.
There are exactly $S_M$ signatures. A stationary work head next sees its
written symbol; an advancing work head sees blank. A permitted input move
preserves the scanned symbol. These facts define an exact finite transition
system, with an exit when the native transition is undefined or attempts a
forbidden direction. The invariant about blank rays is preserved by writes.

The formal simulation lifts every nonexit finite transition to the actual
machine step. Consequently eventual native halting forces eventual exit
from the finite system. Before the first exit its signatures are distinct:
a repeated signature would repeat the future trajectory and prevent exit.
The interval length is therefore bounded by the number of signatures. The
output is reproduced by the native step but need not be in the signature,
since no transition reads it.
\end{outline}

\begin{lemma}[One-pass envelope]\label{lem:envelope}
\formal{RunLength.lean}{pass_envelope}
Suppose the run halts at $T$, $0\le a\le e\le T$, and $[a,e)$ is a pass.
For $L=|w|$,
\[
 L+e+1\le C_M(L+a+1).
\]
\end{lemma}
\begin{outline}
Put $R=L+a$. Every cell visited before $a$, and every nonblank input cell,
lies in $[-R,R]$. Mark a step if some head moves and its departing or
arriving cell lies in this interval. A monotone head has at most $2R+2$
marked steps, so the total number $m$ is at most $(k+1)(2R+2)$.

Between marked steps, a head inside the envelope remains stationary.
A head outside it can move only outward, since it started the pass inside
and is monotone. Its forward work cells are blank: they were visited neither
before $a$ nor earlier in the pass. The input head outside the envelope sees
only the outside-word symbol. Thus each unmarked interval satisfies
\cref{lem:ray}. The $m$ marked steps split $[a,e)$ into at most $m+1$
unmarked intervals, giving $e-a\le(m+1)S_M+m$. The displayed bound follows
from \cref{eq:machine-constants}.
\end{outline}

\begin{theorem}[Halting run-length bound]\label{thm:run-length}
\formal{RunLength.lean}{halting_run_length, halting_time_le, timeBound_log}
If the actual run on an $L$-bit word halts at $T$ and has a partition into
$r$ passes, then
\begin{equation}\label{eq:run-length}
 L+T+1\le (L+1)C_M^r.
\end{equation}
In particular, with $\overline T=(L+1)C_M^r$,
\begin{equation}\label{eq:time-log}
 T\le\overline T,\qquad
 \lceil\log_2(2\overline T+1)\rceil
 \le 2+\lceil\log_2(L+1)\rceil+r\ell_M.
\end{equation}
\end{theorem}
\begin{outline}
Apply \cref{lem:envelope} to successive partition intervals, starting with
$L+0+1$. This gives $L+p_i+1\le(L+1)C_M^i$ at every endpoint $p_i$.
At $p_r=T$ it proves \cref{eq:run-length}. For \cref{eq:time-log}, bound
$L+1$ and $C_M$ by their respective powers of two and use
$2\overline T+1\le4\overline T$, since $\overline T\ge1$.
\end{outline}
This is a restriction on how much computation a halting construction can
hide in a small number of passes. It is not a lower bound on time itself.

\section{Removing time from the adaptive pass bound}\label{sec:adaptive}
Write $H_r=\sum_{i=0}^{r-1}(k+1)^i$, so $kH_r+1=(k+1)^r$, and let
$G_\gamma=\lceil d/\gamma\rceil$ for $\gamma\ge1$. Set
\[
 B_M(L,r)=\sigma_M+(k+4)
   \bigl(2+\lceil\log_2(L+1)\rceil+r\ell_M\bigr).
\]

\begin{theorem}[Time-free grouped bound]\label{thm:grouped}
\formal{AdaptivePasses.lean}{time_free_grouped, time_free_grouped_on, time_free_div}
Suppose $M$ correctly transposes every $d\times d$ input with $b$-bit entries,
and every input's halting run has a partition into $r$ passes. Then
\begin{equation}\label{eq:grouped}
 d^2b\le H_r\bigl((\gamma+1)db+G_\gamma B_M(d^2b,r)\bigr).
\end{equation}
In particular $db\le H_r(2b+B_M(d^2b,r))$.

More generally, let $W$ be a set of inputs on each of which the machine
halts with the correct transposed output and an $r$-pass partition. If
$2^{d^2b}\le2^m|W|$, then
\begin{equation}\label{eq:grouped-on}
 d^2b\le G_\gamma m+
 H_r\bigl((\gamma+1)db+G_\gamma B_M(d^2b,r)\bigr).
\end{equation}
No condition is imposed on inputs outside $W$.
\end{theorem}
\begin{outline}
\Cref{thm:run-length} supplies the common bound
$\overline T=(d^2b+1)C_M^r$ on all relevant halting times. Apply the existing
Lean-verified grouped Theorem P$'$ (or its correct-set form) from the first
paper, and bound its term $\lceil\log_2(2\overline T+1)\rceil$ using
\cref{eq:time-log}. The correct-set form keeps the original payment
$G_\gamma m$. The per-$d$ form uses the original theorem's corresponding
arithmetic consequence.
\end{outline}

\begin{theorem}[Logarithmic passes without a runtime assumption]\label{thm:passes}
\formal{AdaptivePasses.lean}{logarithmic_pass_power, canonical_pass_power, canonical_logarithmic_passes}
Let $k\ge1$, $b\ge1$, and $d\le2^b$. Under the full-correctness and halting
hypotheses of \cref{thm:grouped},
\begin{equation}\label{eq:power}
 kd\le A_M(k+1)^r.
\end{equation}
For the canonical width $b=b(d)$ and $d\ge1$, this gives
\[
 r\ge\log_{k+1}d-
   \bigl(\log_{k+1}A_M-\log_{k+1}k\bigr).
\]
The subtracted quantity depends only on the fixed machine.
\end{theorem}
\begin{outline}
Because $b<2^b$ and $d\le2^b$, one has
$d^2b+1\le2^{3b}$. If $r\le b$, then
$2b+B_M(d^2b,r)\le A_M b$. The per-$d$ bound gives $d\le A_M H_r$;
multiply by $k$ and use $kH_r+1=(k+1)^r$.
If $r>b$, then $d\le2^b\le(k+1)^r$, and $A_M\ge k$ proves the same
conclusion. The canonical width satisfies $b(d)\ge1$ and $d\le2^{b(d)}$.
Taking logarithms in \cref{eq:power} proves the final inequality.
\end{outline}

\begin{corollary}[The earlier runtime-restricted consequence]\label{cor:old}
\formal{AdaptivePasses.lean}{bounded_runtime_corollary}
The conclusion \cref{eq:power} at canonical width still holds if the
hypotheses additionally give a common runtime bound $T_0$. In particular,
the earlier logarithmic-pass consequence for polynomial-time machines
follows without using its polynomial bound.
\end{corollary}
\begin{outline}
Discard the bound $T\le T_0$ and apply \cref{thm:passes}. This rederives the
runtime-restricted consequence, not the original finite-time inequality,
which is an input to our proof.
\end{outline}

\section{Consequences and remaining problem}\label{sec:roadmap}
An adaptive transposer cannot evade the logarithmic-pass barrier by using
very long monotone stretches between direction changes. The grouped
correct-set inequality also removes the need to assume a common runtime
on the specified halting inputs. Its co-entropy cost is retained explicitly;
we do not assert an additional robust asymptotic bound here.

These results do not decide (T), the existence of a fixed finite-tape
machine transposing the canonical inputs in $O(d^2b(d))$ time. A linear-time
run can still have many short passes. A stronger impossibility argument
must turn the required passes into enough total work, while a construction
must exploit them within its time budget. The finite-time P$'$ theorem
also applies to certain nonhalting prefixes; our new run-length argument
requires eventual halting and does not replace that broader interface.
